% ===================== CHAPITRE 4 — SPRINT 2 =========================
\chapter{Sprint 2 : Gestion des agences et des employés}

\section*{Introduction}
Ce sprint met en place les fonctionnalités de gestion du référentiel : les
agences, les employés et leur rattachement. Il suit la même démarche : objectifs,
spécification (cas d'utilisation et backlog), analyse, conception et réalisation.

\section{Objectifs du Sprint}
L'objectif est d'offrir la \textbf{gestion complète (CRUD)} des agences et des
employés, ainsi que le \textbf{rattachement} de chaque employé à son agence et la
gestion des gestionnaires.

\section{Spécification des besoins}

\subsection{Diagramme de cas d'utilisation du Sprint 2}
La figure~\ref{fig:uc_s2} présente les cas d'utilisation du sprint.

\begin{figure}[htbp]
  \centering
  \includegraphics[width=0.82\textwidth]{uc_sprint2.png}
  \caption{Diagramme de cas d'utilisation du Sprint 2}
  \label{fig:uc_s2}
\end{figure}

\subsection{Sprint Backlog du sprint 2}
Le tableau~\ref{tab:bl_s2} détaille les user stories et les tâches techniques.

\begin{table}[htbp]
\centering
\caption{Sprint Backlog du sprint 2}
\label{tab:bl_s2}
\renewcommand{\arraystretch}{1.3}
\begin{tabularx}{\textwidth}{>{\bfseries}p{1cm}p{5.2cm}X}
\toprule
\textmd{\textbf{US}} & \textbf{User story} & \textbf{Tâches techniques}\\
\midrule
US3 & Gérer les agences (créer, consulter, modifier, supprimer). &
Développer l'\texttt{AgenceServlet} et l'\texttt{AgenceService} (EJB) sur l'entité
\texttt{Agence} ; réaliser la page \texttt{agences.jsp} avec les actions CRUD.\\
\midrule
US4 & Gérer les employés et leur rattachement aux agences. &
Développer l'\texttt{AjouterEmployeServlet} et le \texttt{BUtilisateurServlet} ;
gérer les employés via l'\texttt{UtilisateurService} ; associer chaque employé à
une \texttt{Agence} (\texttt{SK\_AGENCE}) ; pages \texttt{employes.jsp} et
\texttt{add-effectif.jsp}.\\
\midrule
US5 & Gérer les gestionnaires et leurs portefeuilles. &
Développer le \texttt{GestionnaireServlet} et le \texttt{GestionnaireService} sur
l'entité \texttt{Gestionnaire} ; page \texttt{gestionnaire.jsp}.\\
\bottomrule
\end{tabularx}
\end{table}

\section{Analyse des besoins}

\subsection{Description textuelle des cas d'utilisation}

\cudesc{tab:cu_agence}{Gérer les agences}
{Permettre à l'administrateur de créer, consulter, modifier et supprimer les
agences du réseau.}
{Administrateur}
{L'administrateur doit être authentifié avec le rôle \texttt{ADMIN}.}
{L'agence créée ou modifiée est enregistrée dans la base et la liste est mise à
jour.}
{(1) L'administrateur accède à l'interface de gestion des agences ; (2) il
sélectionne une action (créer, modifier, supprimer) ; (3) il saisit ou met à jour
les informations de l'agence ; (4) l'\texttt{AgenceService} persiste la
modification et la liste est rafraîchie.}

\cudesc{tab:cu_employe}{Gérer les employés}
{Permettre à l'administrateur de gérer les employés et de les rattacher à leur
agence.}
{Administrateur}
{L'administrateur doit être authentifié et l'agence de rattachement doit exister.}
{L'employé est enregistré et rattaché à son agence.}
{(1) L'administrateur accède à l'interface de gestion des employés ; (2) il saisit
les informations de l'employé ; (3) il sélectionne l'agence de rattachement ; (4)
l'\texttt{UtilisateurService} enregistre l'employé associé à l'agence.}

\section{Conception}

\subsection{Diagramme de séquence : « Créer une agence »}
La figure~\ref{fig:seq_agence} décrit le scénario de création d'une agence.

\begin{figure}[htbp]
  \centering
  \includegraphics[width=0.92\textwidth]{seq_agence.png}
  \caption{Diagramme de séquence -- Création d'une agence}
  \label{fig:seq_agence}
\end{figure}

L'\texttt{AgenceServlet} reçoit la requête de création et délègue à
l'\texttt{AgenceService}, qui persiste la nouvelle \texttt{Agence} via
l'\texttt{EntityManager} (\texttt{INSERT}), puis recharge la liste des agences
renvoyée à la page \texttt{agences.jsp}. Le même schéma s'applique aux opérations
de modification et de suppression.

\section{Réalisation : interfaces utilisateur}
La figure~\ref{fig:ui_agences} illustre les interfaces de gestion des agences et
des employés. \emph{Les captures d'écran réelles seront insérées à cet emplacement.}

\begin{figure}[htbp]
  \centering
  \fbox{\parbox[c][3.6cm][c]{0.8\textwidth}{\centering\itshape Capture à insérer
  -- gestion des agences et des employés}}
  \caption{Interfaces de gestion des agences et des employés}
  \label{fig:ui_agences}
\end{figure}

\section*{Conclusion}
Ce sprint a doté l'application de la gestion du référentiel des agences et des
employés. Le sprint suivant traite des clients et des objectifs commerciaux.
